\documentclass[11pt]{article}
\usepackage{fontspec}
\usepackage[french,english]{babel}
\usepackage{amsmath}
\usepackage{amssymb}
\usepackage{amsthm}
\usepackage{mathrsfs}
\usepackage{tikz-cd}
\usepackage{enumitem}

\newcommand{\printedpage}[1]{\marginpar{\footnotesize p.~#1}}
\newcommand{\msp}[1]{\marginpar{\footnotesize [#1]}}
\newcommand{\defeq}{\overset{\textup{def}}{=}}
\newcommand{\cocart}{\textup{cocart.}}

\begin{document}

\setcounter{footnote}{4}

\printedpage{4}
% head: A. Grothendieck : Dérivateurs, Ch. XVII, Édition des Universités de Montpellier II et Paris VII
% foot: Version 9 février 2019

\msp{page 7}

\[
\begin{tikzcd}
X'_{i_0} \arrow[r] \arrow[d, "\alpha_{i_0}"'] \arrow[dr, phantom, "\cocart"] & X_{i_0} \arrow[r] \arrow[d, "\beta_{i_0}"'] \arrow[dr, phantom, "\cocart"] & X_{\alpha} \arrow[d] \\
Y'_{i_0} \arrow[r] & \tilde{Y}_{i_0} \arrow[r] \arrow[d, "\tilde{u}_{i_0}"'] \arrow[dr, phantom, "\cocart"] & Z'_{i_0} \arrow[d, "v_{i_0}"] \\
& Y_{i_0} \arrow[r] & Z_{i_0}\ .
\end{tikzcd}
\]

We have $\tilde{u}_{i_0} \in C$ by hypothesis, hence $v_{i_0} \in C$, which completes the proof.

\textbf{Proposition 2.} \textit{$(\mathcal{M}, C)$ and where applicable $W$, as in proposition 1 and its corollary, $I$ a locally finite ordered set, i.e. such that $\forall\ i_0 \in I$, $I/i_0$ is finite. Let}
\[
u_{\bullet}\ : X_{\bullet} \longrightarrow Y_{\bullet}
\]
\textit{an arrow in\footnote{in fact, the hypothesis that $X_{\bullet}$, $Y_{\bullet}$ in $\mathcal{M}(I^{\circ})_{\mathrm{c}}$ is unnecessary, both in proposition 2 and in its two corollaries.} $\mathcal{M}(I^{\circ})_{\mathrm{c}} =$ full subcategory of direct systems on $I$, with values in $\mathcal{M}$, such that $\underbrace{\emptyset_{\mathcal{M}(I^{\circ})}}_{\substack{\textup{initial}\\ \textup{object}\\ \textup{of } \mathcal{M}(I^{\circ})}} \longrightarrow Z_{\bullet}$ is in $C_I$ (defined in prop. 1), i.e. such that $\forall\ i_0 \in I$, the arrow}
\[
Z'_{i_0} \ = \ \varinjlim_{i < i_0} Z_i \ \longrightarrow \ Z_{i_0}
\]
\textit{is in $C$, and suppose $u_{\bullet} \in C_I$. Then}
\begin{enumerate}
\item[a)] $\varinjlim\limits_{i \in I} u_{\bullet} : X_I \defeq \varinjlim\limits_{i \in I} X_i \longrightarrow Y_I \defeq \varinjlim\limits_{i \in I} Y_i$ \textit{is in $C$,}
\msp{page 8}\\
\textit{and}
\item[b)] \textit{$\forall\ i_0 \in I$, we have $\alpha_{i_0}, \beta_{i_0} \in C$ in diagram $(*)$ p. 1.}
\end{enumerate}

\textbf{Corollary.} \textit{Moreover (assuming $W$ given as in prop. 1, cor. 1), the following conditions on a $u_{\bullet}$ between $C$-cofibrant $X_{\bullet}$, $Y_{\bullet}$ are equivalent.}
\begin{enumerate}
\item[$1^{\circ}$)] \textit{$u_{\bullet} \in (TC)_I$, i.e. the $\tilde{u}_{i_0} \in TC$ ($\forall\ i_0 \in I$).}
\item[$2^{\circ}$)] \textit{$u_{\bullet} \in T_I \cap W_{I^{\circ}}$, i.e. $\forall\ i_0 \in I$, $\tilde{u}_{i_0} \in C$, $u_{i_0} \in W$.} % [sic] T_I -> C_I, the explanation that follows gives \tilde u_{i_0} \in C, i.e. the condition C_I
\item[$3^{\circ}$)] \textit{For every $i_0 \in I$, we have}
\[
\alpha_{i_0},\ \beta_{i_0},\ \tilde{u}_{i_0},\ u_{i_0} \ \in \ TC
\]
\textit{(i.e. $\alpha_{i_0}, \tilde{u}_{i_0} \in TC$, which implies $\beta_{i_0} \in TC$, then $u_{i_0} \in TC$ by cobase change and by composition).}
\end{enumerate}
\textit{And these conditions imply}
\begin{enumerate}
\item[$4^{\circ}$)] \textit{that $u_I : X_I \longrightarrow Y_I$ is in $TC$.}
\end{enumerate}

Let us see how or to what extent the corollary follows from the proposition. We have of course % [sic] sur -> sûr, accent missing (printed "bien sûr" elsewhere)
\printedpage{5}
% head: A. Grothendieck : Dérivateurs, Ch. XVII, Édition des Universités de Montpellier II et Paris VII
% foot: Version 9 février 2019
\[
\begin{tikzcd}[column sep=small]
& 3^{\circ} \arrow[dl, Rightarrow] \arrow[dr, Rightarrow] & \\
1^{\circ} & & 2^{\circ}
\end{tikzcd}
\]
tautologically. Besides, the proposition applied to $TC$ instead of $X$ shows that we have $1^{\circ} \Longrightarrow 3^{\circ}$ \& $4^{\circ}$, for by applying it to the sets $I_{<i_0}$ we find that $\alpha_{i_0} \in TC\ \forall\ i_0 \in I$. So proposition 2 implies % [sic] X -> C, the proposition is applied to TC instead of C
\[
(1^{\circ} \Longleftrightarrow 3^{\circ}) \Longrightarrow 2^{\circ}\ (\&\ 4^{\circ}).
\]

\msp{page 9}

So to prove $2^{\circ} \Longrightarrow 3^{\circ}$, it would remain to prove that $2^{\circ}$ suffices to imply $4^{\circ}$ (and this for any locally finite ordered set $I$) -- for applying this to the $I_{<i_0}$, we find that $2^{\circ}$ implies $\alpha_{i_0} \in TC = C \cap W$ (not only $\alpha_{i_0} \in C$, according to proposition 2), hence $\beta_{i_0} \in TC \subset W$, hence $\tilde{u}_{i_0} \in W$ by saturation, hence $\tilde{u}_{i_0} \in TC$, whence $3^{\circ}$.

Note on the other hand that b) in proposition 2 is a consequence of a), applied to the restriction of $u_{\bullet}$ to $I_{<i_0}$. All in all, there are therefore only the two things to prove
\[
\left\{
\begin{aligned}
&\text{a)} & u_{\bullet} \in C_I &\quad\Longrightarrow\quad u_I \in C \\
&\text{b)} & u_{\bullet} \in C_I \cap W &\quad\Longrightarrow\quad u_I \in C \cap W\ \ (\defeq TC) \\
& & &\qquad\qquad\ \ (\text{and } u_{\bullet} \in (TC)_I)
\end{aligned}
\right.
\tag*{$(*)$}
\]

We shall prove this by reduction to the special case of proposition 1, already proved.

Besides, I also feel like proving the

\textbf{Corollary 2.} \textit{$\mathcal{M}$, $C$ (and possibly $W$). The hypotheses on $X_{\bullet}, Y_{\bullet} \in \mathcal{M}(I^{\circ})_{\mathrm{c}}$ and on $u_{\bullet} \in C_I$, finally on $I$ locally finite, are those of proposition 2. For every subset $J$ of $I$, let}
\[
u_J \defeq \varinjlim_{i \in J} u_i \ :\ X_J \defeq \varinjlim_{i \in J} X_i \ \longrightarrow\ Y_J \defeq \varinjlim_{i \in J} Y_i\ .
\]
\textit{Let $I' \subset I''$ be two subsets} open \textit{in $I$, and consider the diagram}

\msp{page 10}

\[
\begin{tikzcd}
X_{I'} \arrow[r, "\tau_X"] \arrow[d, "\alpha = u_{I'}"'] \arrow[dr, phantom, "\cocart"] & X_{I''} \arrow[d, "\beta"'] \arrow[dd, bend left=30, "u_{I''}"] & \\
Y_{I'} \arrow[r, "\tilde{\tau}_Y"] \arrow[dr, dotted, "\tau_Y"'] & \tilde{Y}_{I''} \arrow[d, "\tilde{u}"'] & \text{\footnotesize (defined as an amalgamated sum)} \\
& Y_{I''}\ . &
\end{tikzcd}
\tag*{$(*)$}
\]

\printedpage{6}
% head: A. Grothendieck : Dérivateurs, Ch. XVII, Édition des Universités de Montpellier II et Paris VII
% foot: Version 9 février 2019

\textit{Then we have}
\[
\left\{
\begin{aligned}
&\textit{$u_{I'}, u_{I''} \in C$ (hence also $\beta \in C$) (i.e. for every $J$ open in $I$,}\\
&\textit{$u_J : X_J \longrightarrow Y_J$ is in $C$)\,\footnotemark, and $\tau_X, \tau_Y \in C$ (hence also $\tilde{\tau}_Y \in C$)}\\
&\textit{(i.e. for $I' \subset I''$ an inclusion of open subsets, and $X_{\bullet}$ $C$-cofibrant in $\mathcal{M}(I^{\circ})$,}\\
&\textit{we have $X_{I'} \longrightarrow X_{I''}$ in $C$)\,\footnotemark, and finally, even}\\
&\qquad\qquad\qquad\qquad\boxed{\tilde{u} \in C.}
\end{aligned}
\right.
\]
\addtocounter{footnote}{-1}\footnotetext{This is of course contained in prop. 2, applied to $u_{\bullet}|_J$.}\stepcounter{footnote}\footnotetext{It is for this assertion $\tau_X \in C$ (whence $\tilde{\tau}_Y, \tau_Y \in C$, provided $\tilde{u} \in C$) that we need the hypothesis $X_{\bullet} \in \mathcal{M}(I^{\circ})_{\mathrm{c}}$, $\tilde{u} \in C$ true without this hypothesis.}

\textbf{N.B.} The case of proposition 2 in the form $(*)$ a), page 9, (i.e. also the fact that the $u_{I'}$, $u_{I''}$, \dots\ are in $C$) is obtained by taking $I'' = I$, $I' = \emptyset$, the diagram becomes
\[
\begin{tikzcd}
\emptyset \arrow[r] \arrow[d, "\mathrm{id}"'] \arrow[dr, phantom, "\cocart"] & X_I \arrow[d, "\mathrm{id}"] \\
\emptyset \arrow[r] & X_I \arrow[d, "\tilde{u} = u_I"] \\
& Y_I\ ,
\end{tikzcd}
\]
and the conclusion is therefore $u_I \in C$\footnote{The case $(Y_{I'} \xrightarrow{\tau_Y} Y_{I''}) \in C$ is obtained by taking $X_{\bullet} = \emptyset_{\bullet}$, we shall have
\[
\begin{tikzcd}[ampersand replacement=\&]
\emptyset \arrow[r, "\tau_X", "\sim"'] \arrow[d] \arrow[dr, phantom, "\cocart"] \& \emptyset \arrow[d] \\
Y_{I'} \arrow[r, "\mathrm{id}"] \arrow[dr, "\tau_Y"'] \& \tilde{Y}_{I''} = Y_{I'} \arrow[d, "\tilde{u} = \tau_Y"] \\
\& Y_{I''}\ .
\end{tikzcd}
\]}. The \textit{definition} of $C_I$, i.e. the condition $u \in C_I$, is recovered by requiring $\tilde{u} \in C$ in $(*)$ above, in the case $I'' = I/i_0$, $I' = I_{<i_0}$. It is remarkable
\msp{page 11}
that this definition of $C_I$ gives a conclusion as strong as $\tilde{u} \in C$ for \textit{every} pair $I' \subset I''$ of open subsets of $I$.

It is in the end in this form that we are going to prove proposition 2 -- it is \textit{this} now that seems to me \textit{the} right statement in this direction. With the supplement, of course, that if moreover $u_{\bullet} \in W_{I^{\circ}}$, i.e. the $u_i$ are in $W$, then $u_{\bullet}$ is even in $(TC)_I$, or again
\printedpage{7}
% head: A. Grothendieck : Dérivateurs, Ch. XVII, Édition des Universités de Montpellier II et Paris VII
% foot: Version 9 février 2019
simply that the $u_{I'}$ \textit{are in} $TC$ (whence it will therefore follow that $\tilde{u}$ in $(*)$, p. 10, is also in $TC$).

We are going to proceed by transfinite construction. We may of course suppose that $I'' = I$. We are nevertheless going to prove first the special case where $I' = \emptyset$, i.e. prove at least that $u_I \in C$. We construct by transfinite induction a strictly increasing sequence of open subsets
\[
I_{\alpha} \ \subset \ I,
\]
ending at
\[
I_{\alpha_0} \ = \ I,
\]
by setting
\[
\left\{
\begin{aligned}
I_0 \ &= \ \emptyset\,, \\
I_{\alpha} \ &= \ \textstyle\bigcup_{\beta < \alpha} I_{\beta} && \text{if $\alpha$ a limit ordinal}\,, \\
I_{\alpha+1} \ &= \ I_{\alpha} \cup (I/s_{\alpha}) && \text{if $I_{\alpha} \neq I$, and choosing $s_{\alpha} \in I \backslash I_{\alpha}$}\,.
\end{aligned}
\right.
\]

\msp{page 12}

(We end up having $I_{\alpha_0} = I$, of course.) We set
\[
\begin{aligned}
X_{\alpha} \ &= \ X_{I_{\alpha}} \ = \ \varinjlim_{i \in I_{\alpha}} X_i \\
Y_{\alpha} \ &= \ Y_{I_{\alpha}} \ = \ \varinjlim_{i \in I_{\alpha}} Y_i \\
u_{\alpha} \ &= \ u_{I_{\alpha}} \ \ :\ \ X_{\alpha} \longrightarrow Y_{\alpha}\ ,
\end{aligned}
\]
homomorphism of transfinite direct systems. We shall have
\[
X_{\alpha_0} \ = \ X_I, \qquad Y_{\alpha_0} \ = \ Y_I, \qquad u_{\alpha_0} \ = \ u_I,
\]
and we must therefore prove that $u_{\alpha_0} \in C$. For this, we apply proposition 1, so it suffices to prove that $(u_{\alpha}) \in C_{I_{\alpha_0+1}}$ (where $I_{\alpha_0+1} = \{\text{ordinals } i \mid i < \alpha_0 + 1, \text{ i.e. } i \leq \alpha_0\}$). So we must prove that the arrows $\tilde{u}_{i_0}$ in $(*)$ p. 1 (where this time the $i$, $i_0$ are ordinals, $i_0 \leq \alpha_0$) are in $C$. I am going to distinguish the two cases where $\alpha = i_0$ is a limit ordinal, which is trivial for then $X_{i_0} \xleftarrow{\ \sim\ } \varinjlim\limits_{i < i_0} X_i$, $Y_{i_0} \xleftarrow{\ \sim\ } \varinjlim\limits_{i < i_0} Y_i$, and $\tilde{u}_{i_0} = \mathrm{id}$, and the case $\alpha$ not a limit, i.e.
\[
\begin{tikzcd}
X_{I_{\alpha-1}} \arrow[r] \arrow[d, "u_{\alpha-1}"'] & X_{I_{\alpha}} \arrow[d, "\beta_{\alpha}"'] \arrow[dd, bend left=30, "u_{\alpha}"] \\
Y_{I_{\alpha-1}} \arrow[r] & \tilde{Y}_{I_{\alpha}} \arrow[d, "\tilde{u}_{\alpha}"'] \\
& Y_{I_{\alpha}}\ ,
\end{tikzcd}
\]
\printedpage{8}
% head: A. Grothendieck : Dérivateurs, Ch. XVII, Édition des Universités de Montpellier II et Paris VII
% foot: Version 9 février 2019
\msp{page 13}
to prove that $\tilde{u}_{\alpha} \in C$. When we suppose in addition $u_0 \in W_{I_0}$, i.e. the $u_i \in W$, we must in addition prove that $\tilde{u}_{\alpha} \in TC = C \cap W$. Proposition 2 and its corollary 1 will follow from this. To prove $\tilde{u}_{\alpha} \in C$ resp. $\tilde{u}_{\alpha} \in TC$, we may suppose $I_{\alpha} = I$ (changing notation if need be). We therefore set
\[
\begin{aligned}
I_{\alpha} \ &= \ I, & I_{\alpha-1} \ &= \ I' \\
I/s_{\alpha-1} \ &= \ I_0, & I_0 \cap I' \ &= \ I'_0\ ,
\end{aligned}
\]
so that we have a commutative diagram of open inclusions
\[
\begin{tikzcd}
I_0 \cap I' = I'_0 \arrow[r, hook] \arrow[d, hook] & I_0 \arrow[d, hook] \\
I' \arrow[r, hook] & I = I_0 \cup I'.
\end{tikzcd}
\]
We shall use only the fact that $I_0$, hence $I'_0$, is finite, i.e. that in the pair $(I' \subset I)$, $I$ can be written in the form $I = I' \cup I_0$, with $I_0$ an open subset that is \textit{finite,} to conclude from this that in the diagram
\[
\begin{tikzcd}
X_{I'} \arrow[r, "\tau^X_{I',I}"] \arrow[d, "u_{I'}"'] \arrow[dr, phantom, "\cocart"] & X_I \arrow[d, "v_{I',I}"'] \arrow[dd, bend left=30, "u_I"] \\
Y_{I'} \arrow[r] \arrow[dr, "\tau^Y_{I',I}"'] & \tilde{Y}_I \arrow[d, "\substack{\tilde{u} =\\ u_{I',I}}"'] \\
& Y_I
\end{tikzcd}
\]
we have $\tilde{u} \in C$ resp. $\tilde{u} \in TC = C \cap W$.

\msp{page 14}

We need general notation, on the one hand for an inclusion $J \subset K$ of open subsets of $I$, giving rise to the diagram
\[
\begin{tikzcd}
X_J \arrow[r, "\tau^X_{J,K}"] \arrow[d, "u_J"'] \arrow[dr, phantom, "\cocart"] & X_K \arrow[d, "\tilde{u}_{J,K}"'] \arrow[dd, dotted, bend left=30, "u_K"] & \\
Y_J \arrow[r] \arrow[dr, dotted, "\tau^Y_{J,K}"'] & Y_{J,K} \arrow[d, "u_{J,K}"'] & \text{\footnotesize (formerly $\tilde{Y}_K$)} \\
& Y_K\ , &
\end{tikzcd}
\tag*{(1)}
\]
and for inclusions of three nested open subsets $J \subset K \subset L$ of $I$,

\end{document}
